% GENERATED FILE. Edit canonical lessons, then run npm run generate:books.
% canonical-source-hash: fnv1a64:281de83126b93392
% canonical-lessons: TE-C400-ennika, TE-C400-abhyarthi, TE-C400-ennuko, TE-C400-pracaram, TE-C400-mejariti

\chapter{Elections}
\label{ch:te-elections}
\hlchaptermodality{400}

\begin{chapteropening}
\textbf{By the end of this chapter:} \emph{I can talk about elections, candidates and campaigns.}

Complete the last lesson of chapter 400: I can talk about elections, candidates and campaigns.
\end{chapteropening}

\section[\texorpdfstring{\te{ఎన్నిక} (ennika)}{ఎన్నిక (ennika)}]{\te{ఎన్నిక} (ennika) — an election}

\label{lesson:TE-C400-ennika}



\begin{quote}
Before the new one: say the Telugu for injustice, then the Telugu for to honour.
\end{quote}

\subsection*{You'll want to know: \te{ఎన్నిక}}
\textbf{\te{ఎన్నిక}} — \emph{ennika} — \textquotedblleft{}an election\textquotedblright{}.

\begin{quote}
\te{వచ్చే} \te{నెలలో} \te{మా} \te{రాష్ట్రంలో} \te{ఎన్నికలు} \te{జరుగుతాయి}.
\end{quote}

\begin{quote}
\emph{vaccē nelalō mā rāṣṭraṁlō ennikalu jarugutāyi.}
\end{quote}

\begin{quote}
Elections will be held in our state next month.
\end{quote}

Plural: \te{ఎన్నికలు}, \emph{ennikalu}.

Mostly plural: \te{ఎన్నికలు}, the elections. It is built on \te{ఎన్ను}, to count, to reckon.

\begin{cousinweb}
\textbf{\te{ఎన్నిక}} (\emph{ennika}) is native Dravidian: DEDR 793 lists Telugu \emph{ennika}, \textquotedblleft{}counting, number, esteem\textquotedblright{}, beside \emph{ennu}, \textquotedblleft{}to count, to reckon, to think\textquotedblright{}; Telugu now uses it for an election, a choosing by count.
\end{cousinweb}

\begin{grammarlens}[title={What you've built}]
One more word for giving reasons and opinions.
\end{grammarlens}

\subsection*{Your turn}
\begin{itemize}
\raggedright
  \item \emph{Say it:} \emph{ennika}
  \item \emph{Say it:} \emph{ennika}, once more
  \item \emph{Say it:} \emph{vaccē nelalō mā rāṣṭraṁlō ennikalu jarugutāyi.} — Elections will be held in our state next month.
  \item \emph{Recall:} say the Telugu for injustice, then the Telugu for to honour, then say \emph{ennika} again
\end{itemize}

\begin{tcolorbox}[breakable,colback=teal!4,colframe=teal!35!black,title={Before you move on}]
What does \te{ఎన్నిక} mean? (\textquotedblleft{}An election\textquotedblright{}.) Say it once more.
\end{tcolorbox}
\section[\texorpdfstring{\te{అభ్యర్థి} (abhyarthi)}{అభ్యర్థి (abhyarthi)}]{\te{అభ్యర్థి} (abhyarthi) — a candidate; an applicant}

\label{lesson:TE-C400-abhyarthi}



\begin{quote}
Before the new one: say the Telugu for to honour, then the Telugu for an election.
\end{quote}

\subsection*{You'll want to know: \te{అభ్యర్థి}}
\textbf{\te{అభ్యర్థి}} — \emph{abhyarthi} — \textquotedblleft{}a candidate; an applicant\textquotedblright{}.

\begin{quote}
\te{ఆ} \te{ఎన్నికలో} \te{ఐదుగురు} \te{అభ్యర్థులు} \te{ఉన్నారు}.
\end{quote}

\begin{quote}
\emph{ā ennikalō aiduguru abhyarthulu unnāru.}
\end{quote}

\begin{quote}
There are five candidates in that election.
\end{quote}

Plural: \te{అభ్యర్థులు}, \emph{abhyarthulu}.

Also someone who applies for a job or a course.

\begin{cousinweb}
\textbf{\te{అభ్యర్థి}} (\emph{abhyarthi}) is borrowed from Sanskrit \emph{abhyarthin}, \textquotedblleft{}asking, requesting\textquotedblright{} (Monier-Williams): a candidate is one who asks for a post.
\end{cousinweb}

\begin{grammarlens}[title={What you've built}]
One more word for giving reasons and opinions.
\end{grammarlens}

\subsection*{Your turn}
\begin{itemize}
\raggedright
  \item \emph{Say it:} \emph{abhyarthi}
  \item \emph{Say it:} \emph{abhyarthi}, once more
  \item \emph{Say it:} \emph{ā ennikalō aiduguru abhyarthulu unnāru.} — There are five candidates in that election.
  \item \emph{Recall:} say the Telugu for to honour, then the Telugu for an election, then say \emph{abhyarthi} again
\end{itemize}

\begin{tcolorbox}[breakable,colback=teal!4,colframe=teal!35!black,title={Before you move on}]
What does \te{అభ్యర్థి} mean? (\textquotedblleft{}A candidate; an applicant\textquotedblright{}.) Say it once more.
\end{tcolorbox}
\section[\texorpdfstring{\te{ఎన్నుకో} (ennukō)}{ఎన్నుకో (ennukō)}]{\te{ఎన్నుకో} (ennukō) — to elect; to vote in}

\label{lesson:TE-C400-ennuko}



\begin{quote}
Before the new one: say the Telugu for an election, then the Telugu for a candidate.
\end{quote}

\subsection*{You'll want to know: \te{ఎన్నుకో}}
\textbf{\te{ఎన్నుకో}} — \emph{ennukō} — \textquotedblleft{}to elect; to vote in\textquotedblright{}.

\begin{quote}
\te{విద్యార్థులు} \te{ఆమెను} \te{తరగతి} \te{నాయకురాలిగా} \te{ఎన్నుకున్నారు}.
\end{quote}

\begin{quote}
\emph{vidyārthulu āmenu taragati nāyakurāligā ennukunnāru.}
\end{quote}

\begin{quote}
The students elected her class leader.
\end{quote}

Past (I …): \te{ఎన్నుకున్నాను}, \emph{ennukunnānu}.

Close to \te{ఎంచుకో}, to choose, which you know; \te{ఎన్నుకో} is choosing by a vote.

\begin{cousinweb}
\textbf{\te{ఎన్నుకో}} (\emph{ennukō}) is native Dravidian: \emph{ennu}, \textquotedblleft{}to count, to reckon, to judge\textquotedblright{} (DEDR 793), with -\te{కో}, from \emph{konu}, \textquotedblleft{}to take\textquotedblright{} (DEDR 2151): to pick out for oneself.
\end{cousinweb}

\begin{grammarlens}[title={What you've built}]
One more word for giving reasons and opinions.
\end{grammarlens}

\subsection*{Your turn}
\begin{itemize}
\raggedright
  \item \emph{Say it:} \emph{ennukō}
  \item \emph{Say it:} \emph{ennukō}, once more
  \item \emph{Say it:} \emph{vidyārthulu āmenu taragati nāyakurāligā ennukunnāru.} — The students elected her class leader.
  \item \emph{Recall:} say the Telugu for an election, then the Telugu for a candidate, then say \emph{ennukō} again
\end{itemize}

\begin{tcolorbox}[breakable,colback=teal!4,colframe=teal!35!black,title={Before you move on}]
What does \te{ఎన్నుకో} mean? (\textquotedblleft{}To elect; to vote in\textquotedblright{}.) Say it once more.
\end{tcolorbox}
\section[\texorpdfstring{\te{ప్రచారం} (pracāraṁ)}{ప్రచారం (pracāraṁ)}]{\te{ప్రచారం} (pracāraṁ) — a campaign; publicity}

\label{lesson:TE-C400-pracaram}



\begin{quote}
Before the new one: say the Telugu for a candidate, then the Telugu for to elect.
\end{quote}

\subsection*{You'll want to know: \te{ప్రచారం}}
\textbf{\te{ప్రచారం}} — \emph{pracāraṁ} — \textquotedblleft{}a campaign; publicity\textquotedblright{}.

\begin{quote}
\te{ఎన్నికల} \te{ప్రచారం} \te{నిన్నటితో} \te{ముగిసింది}.
\end{quote}

\begin{quote}
\emph{ennikala pracāraṁ ninnaṭitō mugisindi.}
\end{quote}

\begin{quote}
The election campaign ended yesterday.
\end{quote}

\te{ప్రచారం} \te{చేయు}, to campaign, to spread the word.

\begin{cousinweb}
\textbf{\te{ప్రచారం}} (\emph{pracāraṁ}) is borrowed from Sanskrit \emph{pracāra}, \textquotedblleft{}roaming; coming forth, manifestation\textquotedblright{} (Monier-Williams); in Telugu it is the spreading of a word or a cause.
\end{cousinweb}

\begin{grammarlens}[title={What you've built}]
One more word for giving reasons and opinions.
\end{grammarlens}

\subsection*{Your turn}
\begin{itemize}
\raggedright
  \item \emph{Say it:} \emph{pracāraṁ}
  \item \emph{Say it:} \emph{pracāraṁ}, once more
  \item \emph{Say it:} \emph{ennikala pracāraṁ ninnaṭitō mugisindi.} — The election campaign ended yesterday.
  \item \emph{Recall:} say the Telugu for a candidate, then the Telugu for to elect, then say \emph{pracāraṁ} again
\end{itemize}

\begin{tcolorbox}[breakable,colback=teal!4,colframe=teal!35!black,title={Before you move on}]
What does \te{ప్రచారం} mean? (\textquotedblleft{}A campaign; publicity\textquotedblright{}.) Say it once more.
\end{tcolorbox}
\section[\texorpdfstring{\te{మెజారిటీ} (mejāriṭī)}{మెజారిటీ (mejāriṭī)}]{\te{మెజారిటీ} (mejāriṭī) — a majority}

\label{lesson:TE-C400-mejariti}



\begin{quote}
Before the new one: say the Telugu for to elect, then the Telugu for a campaign.
\end{quote}

\subsection*{You'll want to know: \te{మెజారిటీ}}
\textbf{\te{మెజారిటీ}} — \emph{mejāriṭī} — \textquotedblleft{}a majority\textquotedblright{}.

\begin{quote}
\te{ఆమె} \te{పెద్ద} \te{మెజారిటీతో} \te{గెలిచింది}.
\end{quote}

\begin{quote}
\emph{āme pedda mejāriṭītō gelicindi.}
\end{quote}

\begin{quote}
She won by a big majority.
\end{quote}

From English; in plain Telugu, \te{ఎక్కువ} \te{మంది}, most people.

\begin{cousinweb}
\textbf{\te{మెజారిటీ}} (\emph{mejāriṭī}) is a loan from English \emph{majority}.
\end{cousinweb}

\begin{grammarlens}[title={What you've built}]
One more word for giving reasons and opinions.
\end{grammarlens}

\subsection*{Your turn}
\begin{itemize}
\raggedright
  \item \emph{Say it:} \emph{mejāriṭī}
  \item \emph{Say it:} \emph{mejāriṭī}, once more
  \item \emph{Say it:} \emph{āme pedda mejāriṭītō gelicindi.} — She won by a big majority.
  \item \emph{Recall:} say the Telugu for to elect, then the Telugu for a campaign, then say \emph{mejāriṭī} again
\end{itemize}

\begin{tcolorbox}[breakable,colback=teal!4,colframe=teal!35!black,title={Before you move on}]
What does \te{మెజారిటీ} mean? (\textquotedblleft{}A majority\textquotedblright{}.) Say it once more.
\end{tcolorbox}
